%! Right Angles Revisited --- Setting Up Orthogonality — Unit 4, Lesson 4.0 — Lesson Plan
\documentclass[10pt]{article}
\usepackage{linalg-article}
\usepackage{linalg-boxes}
\usepackage{graphicx}   % -article does NOT load graphicx; needed for thumbnails/screenshots
\graphicspath{{images/}}
\newcommand{\TallMath}[1]{$\displaystyle #1\rule[-1.4em]{0pt}{3.2em}$}
\newcommand{\vv}{\mathbf{v}}
\newcommand{\ww}{\mathbf{w}}
\newcommand{\avec}{\mathbf{a}}
\newcommand{\bb}{\mathbf{b}}
\newcommand{\zero}{\mathbf{0}}
\newcommand{\T}{\mathsf{T}}
\newcommand{\UnitNumberName}{Unit 4: Orthogonality \quad}
\newcommand{\LessonNumberName}{Lesson 4.0: Right Angles Revisited --- Setting Up Orthogonality}

\begin{document}
\begin{center}
    {\Large\bfseries \CourseName: \SchoolYear \\
    {\small\bfseries \UnitNumberName \LessonNumberName}}
\end{center}

\begin{tcolorbox}[colback=blush, colframe=burgundy, boxrule=0.9pt, arc=2mm,
    left=2mm, right=2mm, top=1.5mm, bottom=1.5mm]
    \textbf{Primary Objective:} Students can compute the \emph{dot product} of two vectors, find a
    vector's \emph{length} and \emph{unit vector}, and use the rule
    \emph{$\vv\cdot\ww=0 \Leftrightarrow \vv$ and $\ww$ are perpendicular} (Lesson~1.2) to test
    for and build \emph{orthogonal} vectors. They can explain why the four fundamental subspaces
    come in \emph{perpendicular pairs} (the §3.6 picture) and describe where Unit~4 is going:
    dropping a right angle to find the \emph{closest point} --- projection, least squares, and
    orthogonal matrices $Q$.
\end{tcolorbox}

\renewcommand{\arraystretch}{1.4}
\begin{skillbox}[Priority Ideas \& Skills]{goldbox}
\begin{tabularx}{\linewidth}{@{}X|X@{}}
\textbf{Priority Ideas \& Skills} & \textbf{Key Understandings (the ``why'')} \\[2pt]
\hline
\vspace{-6pt}
\begin{itemize}[leftmargin=1.1em, nosep, after=\vspace{2pt}]
  \item Compute a dot product $\vv\cdot\ww$ and read its \emph{sign} (acute / right / obtuse).
  \item Find a length $\|\vv\|=\sqrt{\vv\cdot\vv}$ and the \textbf{unit vector} $\vv/\|\vv\|$.
  \item Test orthogonality: $\vv\perp\ww \Leftrightarrow \vv\cdot\ww=0$; solve for a missing entry.
  \item See the four subspaces as \emph{perpendicular pairs}; preview projection \& least squares.
\end{itemize}
&
\vspace{-6pt}
\begin{itemize}[leftmargin=1.1em, nosep, after=\vspace{2pt}]
  \item The dot product measures alignment; $0$ means the arrows make a \emph{right angle}.
  \item Orthogonality is the one idea the whole unit rests on --- it is Lesson~1.2, reused.
  \item ``Perpendicular'' turns a picture (a right angle) into an \emph{equation} ($\cdot=0$).
  \item The shortest way from a point to a line/plane is along a \emph{perpendicular} --- projection.
\end{itemize}
\end{tabularx}
\end{skillbox}

\begin{skillbox}[{Vocabulary, Concepts, \& Theorems}]{greenbox}
\renewcommand{\arraystretch}{1.3}
\begin{tabularx}{\linewidth}{@{}l X@{}}
\textbf{Dot product $\vv\cdot\ww$} (recall) & $v_1w_1+v_2w_2+\cdots+v_nw_n$: multiply matching components and add (Lesson~1.2). \\
\textbf{Length / norm $\|\vv\|$} (recall) & $\|\vv\|=\sqrt{\vv\cdot\vv}$; the length of the arrow (Lessons~1.0, 1.2). \\
\textbf{Unit vector} (recall) & $\vv/\|\vv\|$: same direction as $\vv$, length $1$. \\
\textbf{Orthogonal / perpendicular} & $\vv\perp\ww$ means the arrows meet at a right angle; the test is $\vv\cdot\ww=0$. \\
\textbf{Orthogonal complement} (recall §3.6) & Two subspaces that are perpendicular and together fill the whole space (e.g.\ $N(A)$ and the row space). \\
\textbf{Projection} (preview 4.2) & The ``shadow'' of a vector on a line or subspace --- the closest point, found by dropping a perpendicular. \\
\textbf{Least squares} (preview 4.3) & The best answer when $A\mathbf x=\mathbf b$ has \emph{no} exact solution: make the error perpendicular to what you can reach. \\
\end{tabularx}
\end{skillbox}

\begin{fixedskillbox}[Activate Prior Knowledge \& Spiral Review]{blush}
    The Warm-Up rehearses the three prerequisite skills this unit is built on, all from Unit~1:
    \textbf{(1)} computing a \emph{dot product} $\vv\cdot\ww$ (Lesson~1.2),
    \textbf{(2)} finding a \emph{length} $\|\vv\|=\sqrt{\vv\cdot\vv}$ and a \emph{unit vector}
    (Lessons~1.0, 1.2), and
    \textbf{(3)} deciding whether two vectors are \emph{perpendicular} by checking
    $\vv\cdot\ww=0$ (Lesson~1.2). Debrief item~3 aloud: ``dot $=0$ means a right angle'' is the
    single fact all of Unit~4 turns on.
\end{fixedskillbox}

\begin{skillbox}[Hook]{blush}
    A hiker stands \emph{off} a straight trail and wants the shortest walk back to it. Sketch the
    trail as a line and the hiker as a point. Ask: \textbf{``Which route back is shortest?''}
    Everyone's instinct is right: the shortest path meets the trail at a \textbf{right angle}. Name
    the payoff: when you cannot land \emph{on} a target exactly --- a line you can't reach, a data
    cloud no straight line passes through --- the \emph{closest} you can get is found by dropping a
    \textbf{perpendicular}, and ``perpendicular'' is just $\vv\cdot\ww=0$. Pose the unit question:
    \textbf{``How do we turn a right angle into an equation we can solve?''} That equation ---
    orthogonality --- runs the whole unit: projection, least squares, and orthogonal matrices.
\end{skillbox}

\begin{skillbox}[Lesson]{blush}
\begin{multicols}{2}
\textbf{1. The dot product, recalled.} From Lesson~1.2, $\vv\cdot\ww=v_1w_1+v_2w_2+\cdots$:
multiply matching components, add. It is one number. Its \textbf{sign} reports the angle: positive
= acute (aligned), $0$ = right angle, negative = obtuse (opposed).

\textbf{2. Length and unit vectors, recalled.} $\|\vv\|=\sqrt{\vv\cdot\vv}$ is the length of the
arrow. Dividing by the length, $\vv/\|\vv\|$, keeps the direction but makes the length $1$ --- a
\textbf{unit vector}. E.g.\ $\vv=\begin{bsmallmatrix} 3\\4 \end{bsmallmatrix}$ has $\|\vv\|=5$ and
unit vector $\begin{bsmallmatrix} 3/5\\4/5 \end{bsmallmatrix}$.

\textbf{3. The one big idea: perpendicular $\Leftrightarrow$ dot $=0$.} Two nonzero vectors meet
at a \textbf{right angle} exactly when $\vv\cdot\ww=0$. This is the whole unit in one line: it
turns the \emph{picture} of a right angle into an \emph{equation}. Check:
$\begin{bsmallmatrix} 3\\4 \end{bsmallmatrix}\cdot\begin{bsmallmatrix} 4\\-3 \end{bsmallmatrix}=12-12=0$
--- perpendicular. And you can \emph{solve} for orthogonality: what $x$ makes
$\begin{bsmallmatrix} 2\\3 \end{bsmallmatrix}\perp\begin{bsmallmatrix} x\\-2 \end{bsmallmatrix}$?
Set $2x-6=0$, so $x=3$.

\columnbreak

\textbf{4. The four subspaces are perpendicular pairs (spiral §3.6).} We closed Unit~3 with the
big picture: for an $m\times n$ matrix, the nullspace $N(A)$ is perpendicular to the row space
$C(A^{\T})$ inside $\mathbb{R}^n$, and $N(A^{\T})\perp C(A)$ inside $\mathbb{R}^m$. \emph{Why}
perpendicular? $A\mathbf x=\zero$ read row by row says (row)$\cdot\mathbf x=0$ --- every row is
orthogonal to $\mathbf x$. That is dot $=0$ again. Unit~4 makes this pairing the main character.

\textbf{5. The road ahead: drop a perpendicular.} When a target $\mathbf b$ is out of reach, the
closest reachable point is the \textbf{projection} --- the foot of the perpendicular, where the
error $\mathbf b-\mathbf p$ is orthogonal to the space. Preview the sequence:
\textbf{4.1} orthogonality of the four subspaces; \textbf{4.2} projections onto lines and
subspaces; \textbf{4.3} least squares (best-fit lines on real data); \textbf{4.4} orthogonal
matrices $Q$ and Gram--Schmidt. Every step is ``make something perpendicular.''
\end{multicols}
\end{skillbox}

\begin{skillbox}[Active Monitoring]{redbox}
Circulate during the notes practice and Tier work. Watch for and correct:
\begin{itemize}[leftmargin=1.2em, nosep]
  \item dot-product slips --- adding the vectors, or multiplying \emph{mismatched} components,
        instead of pairing $v_1w_1+v_2w_2$;
  \item forgetting the square root in $\|\vv\|=\sqrt{\vv\cdot\vv}$, or not dividing by the length
        when asked for a \emph{unit} vector;
  \item reading a \emph{negative} dot product as ``perpendicular'' --- only \emph{zero} is a right
        angle (negative = obtuse).
\end{itemize}
Cold-call prompts: ``What does the \emph{sign} of the dot product tell you?'' ``How do you shrink a
vector to length~1?'' ``Give me a vector perpendicular to $\begin{bsmallmatrix} 2\\5 \end{bsmallmatrix}$
--- how do you know?''
\end{skillbox}

\begin{skillbox}[Group Work \& Differentiation]{redbox}
\textbf{Right Angles, Everywhere} activity (tiered; mirrors the notes).
\begin{multicols}{3}
\textbf{Tier R --- Remediate.} Compute dot products and lengths; find a unit vector; decide
perpendicular / not from the sign of the dot product.

\columnbreak
\textbf{Tier A --- Approaching.} Solve for a missing entry that makes two vectors perpendicular;
check that a matrix's rows are orthogonal to a nullspace vector (the §3.6 pairing on a $2\times2$).

\columnbreak
\textbf{Tier E --- Extension.} \emph{First look at projection:} given a target $\mathbf b$ and a
line direction $\mathbf a$, find the multiple $t\mathbf a$ so the error $\mathbf b-t\mathbf a$ is
perpendicular to $\mathbf a$ --- deriving $t=\frac{\mathbf a\cdot\mathbf b}{\mathbf a\cdot\mathbf a}$.
\end{multicols}
\end{skillbox}

\begin{skillbox}[Individual Work \& Assessment]{redbox}
\textbf{Exit Ticket} (independent, no notes): compute a dot product and decide whether the pair is
perpendicular; find a length and a unit vector; and a \textbf{conceptual/justification check} ---
solve for the entry that makes two vectors orthogonal and explain what $\vv\cdot\ww=0$ means about
the angle. Collect and sort into ``got it / arithmetic slip / concept slip'' to decide whether 4.1
opens with a quick dot-product re-warm.
\end{skillbox}

\begin{skillbox}[Reinforcement \& Extension]{goldbox}
\textbf{Homework:} dot products and their signs; lengths and unit vectors; orthogonality tests and
solve-for-the-entry; one item checking a matrix's rows against its nullspace vector (§3.6); and a
short justification item. \textbf{Extension:} derive the projection multiplier
$t=\frac{\mathbf a\cdot\mathbf b}{\mathbf a\cdot\mathbf a}$ by forcing the error perpendicular --- a
first taste of Lesson~4.2. \textbf{Preview:} Lesson~4.1, \emph{Orthogonality of the Four Subspaces}
--- \emph{why} each pair is perpendicular, and how the two pairs fill their spaces.
\end{skillbox}
\end{document}
